\chapter{ТЕХНИКО-ЭКОНОМИЧЕСКОЕ ОБОСНОВАНИЕ ПРОЕКТА}

\section{Оценка научно-технического уровня разработки}

Разрабатываемая система мониторинга параметров окружающей среды относится к категории систем промышленной автоматизации. В сравнении с существующими аналогами (например, решениями от Xiaomi или Sonoff), проектируемое устройство обладает повышенной надежностью за счет использования протокола MQTT и возможности локальной обработки данных. 

Экономическая целесообразность разработки обусловлена снижением затрат на закупку импортного оборудования и возможностью глубокой интеграции в существующую ИТ-инфраструктуру предприятия.

\section{Расчет затрат на разработку программно-аппаратного комплекса}

Затраты на разработку ($C_{dev}$) включают в себя расходы на оплату труда, отчисления в фонд социальной защиты населения (ФСЗН), накладные расходы и затраты на материалы.

Смета затрат представлена в таблице 5.1.

\begin{table}[h!]
    \centering
    \caption{Смета затрат на разработку системы}
    \begin{tabular}{|l|c|}
        \hline
        Наименование статей затрат & Сумма, руб. \\ \hline
        Материалы и покупные изделия & 450,00 \\ \hline
        Основная заработная плата разработчиков & 2800,00 \\ \hline
        Дополнительная заработная плата (10\%) & 280,00 \\ \hline
        Отчисления в ФСЗН (34\%) & 1047,20 \\ \hline
        Накладные расходы (40\% от основной ЗП) & 1120,00 \\ \hline
        \textbf{Итого себестоимость разработки} & \textbf{5697,20} \\ \hline
    \end{tabular}
\end{table}

\section{Расчет эксплуатационных затрат и экономического эффекта}

Годовые эксплуатационные затраты ($C_{exp}$) определяются по формуле \eqref{eq:eco_exp}:

\begin{equation}
    C_{exp} = C_{el} + C_{maint} + C_{depr}
    \label{eq:eco_exp}
\end{equation}

\begin{formula_where}
    $C_{el}$ --- затраты на электроэнергию, руб.; \\
    $C_{maint}$ --- затраты на техническое обслуживание, руб.; \\
    $C_{depr}$ --- амортизационные отчисления, руб.
\end{formula_where}

Прогнозируемая экономия достигается за счет предотвращения аварийных ситуаций в серверной, которые могут привести к простою оборудования. Условный годовой экономический эффект ($E_{year}$) рассчитывается как:

\begin{equation}
    E_{year} = (P_{loss} \cdot P_{prob}) - C_{exp}
\end{equation}

\begin{formula_where}
    $P_{loss}$ --- потенциальный ущерб от одного часа простоя, руб.; \\
    $P_{prob}$ --- вероятность предотвращения аварии в год.
\end{formula_where}

\section{Расчет срока окупаемости}

Срок окупаемости проекта ($T_{payback}$) в годах определяется отношением капитальных вложений к годовой прибыли (эффекту):

\begin{equation}
    T_{payback} = \frac{C_{dev}}{E_{year}}
\end{equation}



При расчетном экономическом эффекте 3200 руб./год и затратах на разработку 5697,20 руб., срок окупаемости составит:
$T_{payback} = 5697,20 / 3200 \approx 1,78$ года.

\newpage